\documentclass[a4paper, 12pt]{article}

\usepackage{utils}

\renewcommand*{\today}{21 November 2025}

\begin{document}

\hotbox{Analyse 3}{CM 15}{\today}

\subsection{Fonctions de signe non défini : convergence absolue}

\begin{lemme}{Critère de Cauchy}{}
    Soit $F: [a, b[, a \in \R, b \in ]a, +\infty]$, $F$ admet une limite finie en $b$ si et seulement si
    pour tout $\varepsilon > 0$, il existe $V(\varepsilon)$ voisinage de $b$ tel que pour tous $x, y \in [a, b] \cap V(\varepsilon)$.
    $$
    |F(y) - F(x)| \leq \varepsilon
    $$
    ici un voisinage de $b$ c'est un intervalle de la forme $[b - \eta, b + \eta]$ si $b \in \R$
    et de la forme $[c, b[$ si $b = +\infty$.
\end{lemme}

\begin{demonstration}
    \begin{enumerate}
        \item Supposons que la limite existe. Notons la $l$ alors par définition
        $$
        \forall \varepsilon > 0, \exists V_b(\varepsilon), \forall x \in [a, b] \cap V_b(\varepsilon), |F(x) - l| \leq \frac{\varepsilon}{2}
        $$
        d'où le résultat grâce à l'inégalité triangulaire.
        $$
        |F(y) - F(x)| \leq |F(y) - l| + |l - F(x)|
        $$
        \item Réciproquement, On utilise la caractérisation séquentielle : on va démontrer qu'il existe $l \in \R$
        tel que pour toute suite de $[a, b], (x_n) \to b$ alors $F(x_n) \to l$.
        Soit une telle suite, soit $\varepsilon > 0$. Comme $x_n \to b$, il existe $N(\varepsilon)$ tel que $n \geq N \implies x_n \in V_b(\varepsilon)$
        et donc $n \geq N$ et $p \in \N$,
        $$
        |F(x_{n+p}) - F(x_n)| \leq \varepsilon
        $$
        ce qui est précisément la définition de $(F(x_n))$ est de Cauchy. Comme dans $\R$ toute suite de Cauchy
        converge il existe $l \in \R$ tel que $F(x_n) \to l$.

        Il reste à démontrer que pour toute suite $(x_n)$ telle que $x_n \to b$ la limite de $(F(x_n))$ est la même.
        Soient $(x_n), (y_n)$ deux suites telles que $x_n \to b$ et $y_n \to b$ et $l_x, l_y$ les limites de $(F(x_n))$ et $(F(y_n))$ respectivement.
        Soit la suite $(z_n)$ la suite définie par $z_n = x_n$ si $n$ est pair et $z_n = y_n$ si $n$ est impair.
        On a alors $z_n \to b$. Soit $l_z$ la limite de $(F(z_n))$. La sous-suite $(z_{2n}) = (x_{2n})$ converge donc à la fois
        vers $l_x$ et vers $l_z$ donc $l_x = l_z$. De même pour la sous-suite $(z_{2n+1})$ on a $l_y = l_z$. Ainsi $l_x = l_y$.
    \end{enumerate}
\end{demonstration}

\begin{proposition}{}{}
    Soit $f: [a, b[ \to \R, a \in \R, b \in ]a, +\infty]$ une application localement intégrable. Alors $\int_a^bf(t)dt$ converge
    si et seulement si
    $$
    \forall \varepsilon > 0, \exists V_b(\varepsilon), \forall x, y \in V_b(\varepsilon) \cap [a, b[, x \leq y |\int_x^yf(t)dt| \leq \varepsilon
    $$
\end{proposition}

\begin{demonstration}
    Il suffit d'appliquer le critère de Cauchy au fonction $F: [a, b[ \to \R, F(x) = \int_a^xf(t)dt$.
\end{demonstration}

\begin{proposition}{}{}
    Soit $f: [a, b[ \to \R, a \in \R, b \in ]a, +\infty]$ localement intégrable. Si $\int_a^bf(x)dx$ converge absolument alors elle converge.
\end{proposition}

\begin{demonstration}
    En application la proposition précédente, on en déduit que si il y a 
    \begin{hotwarn}
    à finir
    \end{hotwarn}
\end{demonstration}

\begin{remarque}
    La réciproque est fausse en général. Par exemple la fonction $f: [1, +\infty[ \to \R, f(t) = \frac{\sin(t)}{t}$
    est intégrable sur $[1, +\infty[$ mais pas absolument intégrable.
\end{remarque}



\end{document}